\chapter{Sistema corrente}
\label{cap:rad-sistema-corrente}

\section{Situazione di partenza}

Tech4All non sostituisce un sistema software preesistente: interviene su una
prassi. Chi non ha dimestichezza con gli strumenti digitali, di fronte a un
ostacolo, ricorre oggi a tre strategie, tutte con lo stesso limite.

\begin{description}[leftmargin=1.6cm,style=nextline]
  \item[Delega] La persona chiede aiuto a un familiare o a un conoscente più
    esperto. L'ostacolo viene superato, ma da qualcun altro: la competenza non
    viene acquisita e il problema si ripresenta identico alla successiva
    occasione.
  \item[Ricerca autonoma] La persona cerca in rete. Il materiale disponibile
    è però pensato per chi possiede già un linguaggio tecnico di base, non è
    organizzato in un percorso e soprattutto non offre alcun modo di
    verificare se ciò che si è letto è stato compreso.
  \item[Ricorso a un servizio] La persona si rivolge a un negozio o a un
    servizio di assistenza. L'operazione viene eseguita da un operatore a
    fronte di un costo, e ancora una volta senza trasferimento di competenza.
\end{description}

Il denominatore comune è che nessuna delle tre strade produce autonomia. In
tutti i casi manca un elemento: un percorso strutturato, di difficoltà
graduata, accompagnato da una verifica che dia alla persona la misura di ciò
che ha effettivamente imparato.

\diagramma{0.92}{ad-sistema-corrente}{Acquisizione di una competenza digitale
in assenza del sistema.}{fig:rad-ad-corrente}

\section{Limiti che il sistema intende superare}

Dall'analisi della prassi corrente discendono i tre limiti che il sistema
proposto affronta, riportati in \cref{tab:rad-limiti}.

\begin{table}[H]
\centering
\begin{tabular}{@{}L{4.6cm} L{9.6cm}@{}}
\toprule
\textbf{Limite} & \textbf{Risposta del sistema proposto} \\
\midrule
Assenza di un percorso &
Un catalogo di tutorial classificati per categoria, consultabile senza
barriere d'ingresso e ricercabile per parola chiave. \\

Assenza di verifica &
Un quiz facoltativo associato a ciascun tutorial, corretto dal sistema, con
esito immediato e soluzioni mostrate a consegna avvenuta. \\

Assenza di riscontro sul progresso &
Il conteggio dei quiz superati e l'assegnazione automatica di badge al
raggiungimento di soglie prestabilite. \\
\bottomrule
\end{tabular}
\caption{Limiti della prassi corrente e risposta del sistema.}
\label{tab:rad-limiti}
\end{table}

Resta fuori dalla portata del sistema il problema dell'accesso al dispositivo
e alla connessione: Tech4All presuppone che l'utente disponga di un browser e
di una connessione a Internet, e si limita ad abbassare la soglia di
competenza necessaria a usarli con profitto.
